\section{The complete normalized calculus and the envelope repair}
\label{env:calculus}

Let $\operatorname{nf}$ be the unique computation normal form established
below. Canonical numerical terms are the normal forms of
\[
 t::=x\mid0\mid S(t)\mid I(n,a;u.h).
\]
Thus every iterator counter has neither zero nor successor as its root;
all its children are canonical. The update variable $u$ is bound only in
$h$. Put $1=S0$. Formula roots are retained:
\[
 A::=(s=t)\mid Q_y A\mid D_{x,y}A.
\]
For every displayed instance, substitution means capture-free substitution
followed by normalization. In particular, an expression
$I(n,a;u.h)$ denotes the normal form of the displayed constructor
expression, even when substitution exposes a computation. Formula
normalization changes numerical subterms and never removes a $Q$ or $D$.

The intended readings, extended recursively to matrices, are
\[
 Q_yB(y)\;:=\;B(0)\to\forall y\,B(y),\qquad
 D_{x,y}B(x,y)\;:=\;\exists x\bigl(B(x,0)\to\forall y\,B(x,y)\bigr).
\]
The designated contradiction is $0=S0$. Contexts are finite reusable lists.

\subsection{Six primitive inference patterns}

The calculus $\mathcal F$ has exactly the following primitive patterns.
\begin{enumerate}
\item \emph{Hypothesis:} $\Gamma\vdash A$ if $A\in\Gamma$.
\item \emph{Reflexivity:} $\Gamma\vdash t=t$ for every canonical term $t$.
\item \emph{Crossed transport in one term:}
\[
 \frac{\Gamma\vdash s=t\qquad\Gamma\vdash h(s)=r}
      {\Gamma\vdash r=h(t)}.
 \tag{Tr}\label{env:tr}
\]
Here $h(v)$ is any fixed numerical term template, including a constant
template. There is no restriction on the source $s$.
\item \emph{Closure elimination:}
\[
 \frac{\Gamma\vdash Q_yB(y)\qquad\Gamma\vdash B(0)}
      {\Gamma\vdash B(t)}.
 \tag{QE}\label{env:qe}
\]
\item \emph{Paired elimination:}
\[
 \frac{\Gamma\vdash D_{x,y}B(x,y)\qquad
       \Gamma,Q_yB(b,y)\vdash A}{\Gamma\vdash A}.
 \tag{DE}\label{env:de}
\]
The eigenvariable $b$ is fresh for the original context, conclusion, and
fixed matrix parameters.
\item \emph{Witness-directed introduction:}
\[
 \frac{\Gamma,B(t,\eta)\vdash B(t,Sz)}
      {\Gamma\vdash D_{x,y}B(x,y)},\qquad
 \eta=\begin{cases}
 z,&B\text{ is an equation and }t\text{ is literally }0,\\
 0,&\text{otherwise}.
 \end{cases}
 \tag{WI}\label{env:wi}
\]
Here $z$ is fresh for the context, witness $t$, and fixed matrix parameters.
\end{enumerate}
These are all the logical and nonlogical inference patterns. There is no
additional computation axiom: both iterator equations are enforced by
normalized substitution. Their implementation cost remains part of the
checker. Ordinary logical connectives are subsequently recovered as in the
native scalar calculus; they are not extra primitive rules.

Weakening, fresh renaming, and ordinary-assumption cut follow directly by
induction on derivations. They preserve the literal-zero and formula-root
tests. Numerical substitution is \emph{not} claimed by that direct argument:
it can normalize a nonzero-spelled witness to zero. It will follow from the
comparison with a substitution-stable calculus below.

\subsection{The compound envelope}

\paragraph{First, general closure introduction by cases.}
For every matrix $B$, the rule
\[
 \frac{\Gamma,B(0)\vdash B(Sz)}{\Gamma\vdash Q_yB(y)}
 \tag{QC}\label{env:qc}
\]
is already admissible. Introduce $D_{x,y}B(y)$ with an unused witness
coordinate $x$ and witness $1$. Since $1$ is not literally zero, WI uses
its case branch regardless of the root of $B$. DE and its local hypothesis
then export $Q_yB(y)$. This construction does not use any omitted atomic
case-introduction instance at witness zero.

\paragraph{An equation has a compound representative.}
For $E$ the equation $s=t$, define
\[
 G(v,E):=\bigl(I(v,0;u.s)=I(v,0;u.t)\bigr),\qquad
 L(E):=Q_vG(v,E),
 \tag{Envelope}\label{env:lift}
\]
where $u,v$ are fresh; the endpoints are lifted under these binders.
Computation gives the literal normalized identities
\[
 G(0,E)\equiv(0=0),\qquad G(Sw,E)\equiv E.
\]
Consequently, under every retained context,
\[
 E\vdash L(E),\qquad L(E)\vdash E.
 \tag{Envelope equivalence}\label{env:equiv}
\]
For the first direction apply QC to the gate matrix: under its local
assumption $0=0$, the successor instance is the retained equation $E$.
For the reverse direction use QE at $1$ and reflexivity $0=0$.
The envelope has a $Q$ root after normalization. These derivations are
uniform in all parameters of $E$ and use neither arithmetic induction
nor a numerical case rule.

\paragraph{Recovering the missing atomic case.}
Suppose a supplied derivation proves
\[
 \pi:\quad\Gamma,E(0,0)\vdash E(0,Sz),
 \qquad E(x,y)\text{ an equation}.
\]
Set $F(x,y)=L(E(x,y))$. Under $F(0,0)$, the reverse envelope direction
supplies $E(0,0)$. Cut with $\pi$ and then use the forward envelope
direction to obtain $F(0,Sz)$. WI at witness zero now uses \emph{cases},
because $F$ has a compound root, and proves $D_{x,y}F(x,y)$.

Eliminate this formula at a fresh $b$, retaining $Q_yF(b,y)$. To introduce
$D_{x,y}E(x,y)$ at witness $b$, assume $E(b,0)$, convert it to $F(b,0)$,
use the retained closure at $Sz$, and convert $F(b,Sz)$ back to
$E(b,Sz)$. The fresh variable $b$ is not literally zero, so this second
WI also uses cases. DE closes the derivation, yielding the required
$D_{x,y}E(x,y)$. Neither the final formula nor the original context contains
$b$.

It follows that unrestricted case introduction is admissible:
\[
 \frac{\Gamma,B(t,0)\vdash B(t,Sz)}
      {\Gamma\vdash D_{x,y}B(x,y)}.
 \tag{D-case}\label{env:dcase}
\]
All instances except the repaired one are already WI. This replaces the
earlier nonliteral-zero-alias argument. Full normalization erases closed
zero aliases; the compound envelope works without one.

\paragraph{Direct equation progress.}
From $\Gamma,E(z)\vdash E(Sz)$, with $E$ an equation, use an unused
witness coordinate and witness zero in WI. DE then exports
\[
 \Gamma\vdash Q_yE(y).
 \tag{Q-progress}\label{env:qprog}
\]
No case repair is needed for this construction.

\subsection{Exact comparison and the interpretations}

Let $\mathcal F^+$ be the substitution-stable auxiliary calculus on canonical
terms, with Hyp, reflexivity, Tr, QE, DE, D-case, and Q-progress. The
preceding constructions derive every rule of $\mathcal F^+$ in
$\mathcal F$. Conversely, a case-mode WI is D-case. For a progress-mode
WI, its equation matrix and witness zero give Q-progress for $E(0,y)$;
under $E(0,0)$, QE supplies $E(0,Sz)$, and D-case at zero gives the same
conclusion. The other rules are common. Thus, for every canonical sequent,
\[
 \Gamma\vdash_{\mathcal F}A
 \quad\Longleftrightarrow\quad
 \Gamma\vdash_{\mathcal F^+}A.
 \tag{Exact consequence}\label{env:consequence}
\]

Ordinary equality reasoning is available from reflexivity and Tr.
Symmetry follows by taking the major premise $s=s$ and minor premise
$s=t$ with the identity template. Congruence follows by taking a
reflexive minor $h(s)=h(s)$ along $s=t$. Two crossed transports give
transport in an arbitrary equation template. Transport in compound
formula templates follows by external induction on their roots: for $Q$,
use QC, transport its base backwards, apply the retained closure, and
transport its successor conclusion forwards; for $D$, eliminate the old
formula and apply D-case at the resulting witness with the same two
conversions. This is a syntactic induction on formulas, not an additional
arithmetic-induction schema.

